\section{Rigid-charge particle--rotor control}\label{app:rigid}

Using standard particle--rotor kinematics \cite{BohrMottelson1975},
we assemble the unpaired Hamiltonian directly:
\begin{align}
H(C)&=\sum_{k=1}^3 A_k R_k^2+C\,1\otimes Q_j(\gamma),\\
R_k&=I_{k,\mathrm{right}}\otimes1-1\otimes j_k,\nonumber\\
A_k&=\frac{\hbar^2}{2\mathcal J_k},\quad
\mathcal J_k=\frac43\mathcal J_0\sin^2(\gamma-2\pi k/3).
\end{align}
\begin{equation}
\begin{aligned}
Q_j(\gamma)={}&\frac{\cos\gamma}{4j(j+1)}[3j_z^2-j(j+1)1]\\
&+\frac{\sqrt3\sin\gamma}{4j(j+1)}(j_x^2-j_y^2).
\end{aligned}
\end{equation}
The angular momentum components are dimensionless, in units of $\hbar$;
$\mathcal J_0=30\,\hbar^2/{\rm MeV}$, $j=11/2$ and
$\gamma=20^\circ$. The unsorted principal inertias are
38.79385, 16.52704 and $4.67911\,\hbar^2/{\rm MeV}$.
The field $Q_j$ is dimensionless; its normalization defines $C$ and
is distinct from the core-only electric tensor.
There is no pairing projection or zero-gap approximation.

The right-body representation is
$I_{\rm right}=(J_x,-J_y,J_z)$, with
$[I_1,I_2]=-iI_3$; the particle matrices have ordinary commutators.
The full physical $D_2$ subspace is retained at each
$I=1/2,\ldots,15/2$. These 216 multiplets exhaust all required
$7/2$--$15/2$ intermediate states for initial $I=11/2$ quadratic
and cubic closure. They are not a globally complete infinite-rotor space:
for an initial $I=15/2$, additional $17/2$ and $19/2$ channels would be required.

Equation~(\ref{eq:rigididentity}) follows from rotation invariance of
the intrinsic tensor polynomials. Equivalently, its real traceless
Cartesian tensor has
$\operatorname{tr}A^2=1$ and
$\sqrt6\,\operatorname{tr}A^3=\cos3\gamma$ for $\mathcal Q=1$.
The numerical contractions equal these constants times the identity
on the 36-dimensional initial-$I=11/2,M=1/2$ coordinate space.
Independent Euler-angle integration agrees with the product-basis
matrix elements within $2.8\times10^{-14}$.
The independent $D_2$ group projector, tensor covariance and Hermiticity
are checked as well.

At $C=0$, the exact $R_{\rm core}=0$ spectator has zero rotational
energy and static moment. Its spin fractions equal
$(2I_f+1)/[5(2j+1)]$. The two core-$R=2$ energies, 0.126795 and
0.473205 MeV, are recovered independently from a five-dimensional
integer-spin rotor. A deliberately incorrect body-axis sign violates
the zero-energy spectator condition, and omission of the $15/2$
outer sector loses $4/15$ of the initial quadratic closure.
These checks test different parts of the construction without fitting
the selected ratio.

For a fixed axial $K=J$ rotor with intrinsic tensor amplitude
$\mathcal Q=T_{20}^{\rm body}$, the complete moments are
$q_2=\mathcal Q^2$ and $q_3=\mathcal Q^3$.
The conventional electric intrinsic quadrupole moment instead satisfies
$\mathcal Q=\sqrt{5/(16\pi)}Q_0$.
This control illustrates the regime dependence of Eq.~(\ref{eq:proxy}):
at $J=11/2$ it gives $10.193^\circ$ despite axial input.
These controls validate the definitions and implementation; they are not new
shape theorems or empirical uncertainty estimates.

\section{Conditional charge calibrations}\label{app:calibrated}

For each distinct central $^{136}$Ba $2_1^+$ moment target, the
boson charge is kept positive and recalibrated to the common
$B_\uparrow=0.407\,e^2{\rm b}^2$. Its moment determines $\chi'$,
and $Q_s(^{137}{\rm Ba})=0.78\,e\,{\rm b}$ then determines $e_F$.
All finite sign-consistent roots are retained; the $^{135}$Ba
moment is not used. Its boson charge matches the unchanged
$^{134}$Ba core strength.

\begin{table}[t]
\caption{All conditional moment-calibration branches at fixed
Hamiltonians. $t_{136}$, $e_F$ and $Q_s(135)$ are in $e\,{\rm b}$.
Ratios refer to the lowest $^{135}$Ba $11/2^-$ state.
Both $|r|>1$ and $z_\infty>1$ imply an undefined real angle.}
\label{tab:calibrated}
\centering
\setlength{\tabcolsep}{3pt}
\begin{tabular}{rrrrrr}
\toprule
$t_{136}$ & $\chi'$ & $e_F$ & $Q_s(135)$ & $r$ & $z_\infty$\\
\midrule
$-0.19$ & 0.982 & 0.199 & 1.089 & $-0.642$ & 1.030\\
$+0.01$ & 1.959 & 0.311 & 0.923 & $-1.093$ & 0.850\\
$+0.02$ & 2.004 & 0.317 & 0.915 & $-1.096$ & 0.840\\
$+0.07$ & 2.222 & 0.344 & 0.873 & $-1.092$ & 0.792\\
$+0.25$ & 2.943 & 0.445 & 0.723 & $-0.920$ & 0.640\\
\bottomrule
\end{tabular}
\end{table}

The independent E2 parameter $\chi'$ is not restricted here by a
customary interval for the Hamiltonian parameter $\chi$.
These are central-value comparisons at a common averaged core
strength: moment and strength branches may covary, and the 135/137Ba
moment evaluations share source lineage. No joint confidence region
or preference for a branch is inferred.

A marginal domain crossing illustrates the separate role of
state-count truncation. For $^{137}$Ba with the
$+0.01\,e\,{\rm b}$ core target, $z_5=1.00027$ but
$z_\infty=0.99484$, so completing that spin channel restores its
real selected angle. The complete ratio remains $r=-1.08957$.
This differs from the primary $^{135}$Ba crossing that persists
after within-spin completion.

\section{Isomer units, mappings and decay completeness}\label{app:isomer}

The experimental parent and daughter are the tentative $(19/2^-)$
2349.9-keV and $(15/2^-)$ 2229.8-keV levels, respectively.
Their common anchor is the 661.7-keV $11/2^-$ isomer.
For every open model branch, the photon rate used is
\begin{equation}
\lambda_{\gamma,f}(E2)=1.2239\times10^{13}
 E_{if}^{5}B(E2;i\to f),
\end{equation}
with energy, strength and rate units MeV, $e^2{\rm b}^2$ and s$^{-1}$.
No unobserved modeled daughter is treated as an experimentally excluded level.

The ordinary downward Weisskopf reference is
\begin{equation}
B_{\rm W}(E2)=\frac1{4\pi}\left(\frac35\right)^2
(1.2A^{1/3})^4\,e^2{\rm fm}^4.
\end{equation}
For $A=137$, it equals $0.00419544833\,e^2{\rm b}^2$.
Thus the reported $0.46(3)$ W.u.\ corresponds to
$(1.930\pm0.126)\times10^{-3}\,e^2{\rm b}^2$,
with only the source's quoted margin converted.
Our downward strength is
$|R_{15/2,19/2}|^2/20$, while the reverse strength is larger
by $20/16$. This is different from the $16/12$ factor in
Eq.~(\ref{eq:staticratio}) and the factor five for a $0\leftrightarrow2$
core transition.

\begin{table}[t]
\caption{All four predeclared state pairs for the primary $^{137}$Ba
operator. Relative energies use the lowest model $11/2$ state as
zero. The experimental candidate energies on the same reference
are 1688.2 and 1568.1 keV, respectively. Model-state ordinals
do not establish experimental identities.}\label{tab:isomerpairs}
\centering
\begin{tabular}{crrr}
\toprule
$19/2_i\to15/2_f$ & $E_i$ (keV) & $E_f$ (keV) & $B_\downarrow$ (W.u.)\\
\midrule
$1\to1$ &2142.484&968.874&22.46885\\
$1\to2$ &2142.484&2094.159&2.17762\\
$2\to1$ &3648.848&968.874&0.05585\\
$2\to2$ &3648.848&2094.159&13.14225\\
\bottomrule
\end{tabular}
\end{table}

All seven operators---the primary, core-only and five conditional
calibrations---were evaluated for every pair.
The energy-distance mapping was fixed before the E2 evaluation.
The nearest of the four pairs is $19/2_1\to15/2_2$, whose parent
and daughter residuals are still +454 and +526 keV. This minimum
distance is not a certified physical correspondence or a statistical
chi-square fit.
The all-open-width test was defined afterward as a separate
follow-up, without refitting either Hamiltonian or operator.
The full package retains 28 strengths, 14 parent/operator sums and
84 open E2 branches.

For the primary operator, the two parent half-life upper bounds
are $2.69817\times10^{-4}$ and $6.02153\times10^{-5}$ ns.
With the conditional $t_{136}=+0.07\,e\,{\rm b}$ operator,
the alternative $1\to2$ strength is 0.42002 W.u.,
outside the quoted interval $[0.43,0.49]$ W.u.
Its modeled photon energy is only 48.325 keV.
The competing $1\to1$ branch has 21.86655 W.u. and accounts for
99.9981\% of that model's E2 photon rate. These are model rate fractions,
not measured total branching probabilities.

No experimental photon energy is substituted into a different
model branch. Missing conversion and non-E2 rates are not assumed
zero to claim a complete lifetime prediction. Their nonnegativity
gives the conditional upper bound in Eq.~(\ref{eq:isomerbound});
it cannot compensate for the fast E2 routes of the unchanged model scheme.
